\documentclass{situationist-preprint}

\title{The Epistocratic Objection Answered}
\msid{POL-04-b}
\msfamily{POL-04}
\msdiscipline{Political Philosophy}
\mstype{Companion manuscript}
\mskeywords{epistocracy; democracy; political ignorance; jury theorem; Brennan; Estlund; plural voting}

\begin{document}
\maketitle

\begin{abstract}
The epistocratic objection holds that since voters are widely ignorant, democracy is unreliable, and that rule by the better informed would produce better outcomes. POL-04-a argues that democracy is justified both intrinsically and instrumentally and that neither justification suffices alone; its reverse jury theorem concedes that majority rule amplifies incompetence. We answer the objection in three steps. Ignorance of the facts epistocrats measure is not incompetence on the questions democracies decide. A reliable epistocracy, granted for argument, is unjustified because it fails the intrinsic requirement, and weighted voting fails it too. And the selection of the knowers is itself a question on which reasonable persons disagree, so that no epistocracy can be justified to all of them.
\end{abstract}

\section{The objection}

\citet{Brennan2016} argues that most citizens are badly informed about politics, that their votes are therefore unreliable guides to good policy, and that a system restricting or weighting votes by competence would produce better outcomes; since no one has a right to impose incompetent decisions on others, epistocracy is preferable where it would work better. The empirical premise draws on a large literature on political ignorance \citep{Somin2013}. The normative premise echoes Mill's proposal of plural votes for the educated \citep[ch.~8]{Mill1861}.

POL-04-a meets the objection in part. Its main theorem (Theorem~4.4 of POL-04-a) holds that democracy is justified because it is both egalitarian and reliable, and that a reliable epistocracy is not justified (Assumption~3.4 of POL-04-a). But its reverse jury theorem (Proposition~4.2 of POL-04-a) concedes the mechanism the epistocrat relies on: when individual competence falls below one half, majority rule amplifies error. If the epistocrat's evidence shows competence below one half, the instrumental justification fails (Corollary~4.3 of POL-04-a), and the intrinsic one cannot justify democracy alone. The objection must therefore be answered, not only resisted.

\section{Three replies}

\emph{First, ignorance is not incompetence on the question.} The competence that the jury theorem requires is competence on the questions put to the vote. Surveys measure knowledge of facts, such as the size of the budget or the name of a legislator, that voters need not possess to answer the questions elections pose: which of the available parties or candidates to entrust with government, and whether the incumbents have governed acceptably. These are questions on which citizens can draw on their own experience of the outcomes of policy, and on the cues of trusted intermediaries. Low factual knowledge is compatible with competence above one half on such questions, and the evidence that would show competence below one half on them is evidence of a different kind. Where the questions are technical, democracies delegate them to agencies under democratic control, which is a division of labor, not an epistocracy.

\emph{Second, reliability does not suffice.} Suppose, for the argument, that an epistocracy would be more reliable. It is still not justified, because a procedure is justified only if it gives citizens equal standing as well as tracking good outcomes. The epistocrat may reply that weighted voting gives everyone some say. But equal standing requires an equal say, not some say; a plural-vote scheme declares some citizens' judgement worth more than others' in the making of laws that bind them all, and that is the denial of equal standing that Assumption~3.4 of POL-04-a records. The structure of the main theorem of POL-04-a is designed for this case: reliability is necessary, and not sufficient.

\emph{Third, the knowers cannot be identified in a way all can accept.} Any epistocracy must select its knowers by some test. \citet{Estlund2008} argues that for any proposed test, reasonable persons can object that it is biased, or that it correlates with characteristics that distort political judgement in ways the test does not detect; and that a claim to rule over others must be acceptable to all qualified points of view. We add that this is an instance of the criterion of POL-02: since the choice of test is itself subject to reasonable disagreement, an epistocracy cannot be justified to all reasonable persons subject to it, whereas a procedure that gives each an equal say needs no such test.

\section{Objections}

\begin{objection}
The first reply proves too much: if voters are competent on the questions elections pose, the instrumental case needs no help from the intrinsic one.
\end{objection}

\begin{reply}
It does not prove that. Competence above one half on the questions put to the vote is a condition that holds in favorable circumstances and that the procedure itself helps to sustain; the reverse jury theorem shows what follows where it fails. The intrinsic justification is needed to explain why democracy, and not a more reliable alternative, is the right response to such failures.
\end{reply}

\begin{objection}
The third reply would rule out any expertise in government, including the agencies the first reply relies on.
\end{objection}

\begin{reply}
Expertise in administration under democratic control does not claim authority to make law over others on the ground of superior judgement; it is answerable to a procedure that gives each citizen an equal say. The objection to epistocracy is to the allocation of the final say, not to the use of knowledge.
\end{reply}

\section{Conclusion}

The epistocratic objection rests on a premise about competence that its evidence does not establish, on a conception of justification as reliability alone, and on a selection of knowers that cannot be justified to all reasonable persons. Each failure confirms the conclusion of POL-04: democracy is justified intrinsically and instrumentally, and neither justification suffices alone.

\begin{thebibliography}{9}
\bibitem[Brennan(2016)]{Brennan2016} Brennan, J. (2016). \emph{Against Democracy}. Princeton University Press.
\bibitem[Estlund(2008)]{Estlund2008} Estlund, D.~M. (2008). \emph{Democratic Authority: A Philosophical Framework}. Princeton University Press.
\bibitem[Mill(1861)]{Mill1861} Mill, J.~S. (1861). \emph{Considerations on Representative Government}. Parker, Son, and Bourn.
\bibitem[Somin(2013)]{Somin2013} Somin, I. (2013). \emph{Democracy and Political Ignorance: Why Smaller Government Is Smarter}. Stanford University Press.
\end{thebibliography}

\end{document}
